\documentclass[12pt]{article}
\usepackage[utf8]{inputenc}
\usepackage[T1]{fontenc}
\usepackage[margin=1in]{geometry}
\usepackage{setspace}
\usepackage{xcolor}
\usepackage[colorlinks=true, urlcolor=blue!60!black, linkcolor=blue!60!black]{hyperref}
\setlength{\parindent}{2em}
\setlength{\parskip}{1.25em}
\renewcommand{\baselinestretch}{1.0}
\begin{document}
\begin{center}
{\LARGE \textit{Bradwell v. Illinois}, 83 U.S. 130 (1873)\par}
\vspace{0.8em}
{\large Case Brief\par}
\end{center}
\vspace{1em}

\section{Memory Jogger}

One day after gutting the Fourteenth Amendment's Privileges or Immunities Clause in the \textit{Slaughter-House Cases}, the Supreme Court held 8--1 that Illinois could bar Myra Bradwell from practicing law because admission to a state bar is not a privilege of national citizenship, while Justice Bradley's notorious concurrence invoked ``the law of the Creator'' to confine women to the domestic sphere.

\section{Detailed Case Facts}

Myra Bradwell, a longtime resident of Chicago, Illinois, and the founder and editor of the influential \textit{Chicago Legal News}, applied to the justices of the Supreme Court of Illinois for a license to practice law. As recorded in the \href{https://www.courtlistener.com/opinion/88662/bradwell-v-state/}{CourtListener opinion} and \href{https://supreme.justia.com/cases/federal/us/83/130/}{Justia full text}, she accompanied her petition with the standard certificate from an inferior court attesting to her good moral character and certifying that, upon due examination, she possessed the requisite legal qualifications. Pending her application, she also filed an affidavit stating that she had been born in Vermont, had been a citizen of Vermont, was a citizen of the United States, and had resided for many years in Chicago. She asserted that under Article IV, \S{} 2 of the U.S. Constitution (the Privileges and Immunities Clause) and Section 1 of the Fourteenth Amendment (the Privileges or Immunities Clause), she was constitutionally entitled to a law license.

The governing Illinois bar-admission statute provided that no person could practice as an attorney or counselor at law in any state court of record without first obtaining a license from two justices of the Illinois Supreme Court. Its only express statutory prerequisite was a county-court certificate of good moral character, leaving other admission rules to the state court's discretion. When Bradwell's application first came before the Illinois Supreme Court, the court refused her a license on the ground of coverture---stating that as a married woman she would be bound neither by express contracts nor by the implied contracts that the law creates between attorney and client. After Bradwell filed a printed argument challenging that ground, the Illinois Supreme Court issued a written opinion declaring that because Illinois had adopted the common law of England---where female attorneys were unknown in Westminster Hall---the legislature had never contemplated extending the office of attorney at law to women.

Historically and politically, \textit{Bradwell v. Illinois} was decided on April 15, 1873, just one day after the \textit{Slaughter-House Cases}, 83 U.S. 36 (1873), as documented by the \href{https://constitutioncenter.org/the-constitution/supreme-court-case-library/bradwell-v-illinois-1873}{National Constitution Center}. Both decisions represented the Supreme Court's earliest interpretations of the Fourteenth Amendment, ratified in 1868 during Reconstruction. Women's rights advocates had hoped that Section 1's broad guarantee of the ``privileges or immunities of citizens of the United States'' would dismantle common-law disabilities imposed on women, including occupational exclusions and coverture. Instead, the Court's narrow construction of national citizenship in \textit{Slaughter-House} foreclosed Bradwell's federal claim. (Illinois subsequently enacted legislation opening the bar to women; as noted in \href{https://ideaexchange.uakron.edu/cgi/viewcontent.cgi?article=1164&context=akronlawreview}{Akron Law Review scholarship}, Bradwell refused to reapply, but the Illinois Supreme Court belatedly admitted her on her original motion in 1890, and she was admitted to the U.S. Supreme Court bar in 1892, two years before her death.)

\section{Procedural History}

\begin{itemize}
  \item \textbf{Illinois Supreme Court:} Myra Bradwell petitioned directly for a license to practice law. The court initially refused on coverture grounds and, after receiving her printed argument, issued a written opinion denying her application on the ground that the state legislature had not intended to authorize women to be admitted to the bar.
  \item \textbf{Basis for Supreme Court Review:} Bradwell sued out a writ of error to the U.S. Supreme Court under \S{} 25 of the Judiciary Act of 1789, as amended by the Act of February 5, 1867, claiming that the Illinois Supreme Court's refusal violated Article IV, \S{} 2 and the Fourteenth Amendment of the federal Constitution. The case was argued on January 18, 1873, and decided on April 15, 1873.
  \item \textbf{Disposition:} Judgment of the Supreme Court of Illinois \textbf{affirmed} by an 8--1 vote.
\end{itemize}

\section{Judicial Votes}

\begin{itemize}
  \item \textbf{Majority (5 Justices):} Justice Samuel Freeman Miller (author), joined by Justices Nathan Clifford, David Davis, William Strong, and Ward Hunt.
  \item \textbf{Concurrence in the Judgment (3 Justices):} Justice Joseph P. Bradley (author), joined by Justices Noah Haynes Swayne and Stephen Johnson Field.
  \item \textbf{Dissent (1 Justice):} Chief Justice Salmon P. Chase (dissented from the judgment and from all opinions without writing an opinion, as he was gravely ill and died weeks later on May 7, 1873).
\end{itemize}

\section{Holding}

\begin{itemize}
  \item \textbf{Judgment:} The Supreme Court affirmed the judgment of the Supreme Court of Illinois denying Myra Bradwell a license to practice law.
  \item \textbf{Rule Announced:} First, Article IV, \S{} 2's Privileges and Immunities Clause applies only to citizens of other states and offers no protection to a citizen against the laws of her own state of residence; under the Fourteenth Amendment's Citizenship Clause, Bradwell's longtime Chicago residence made her a citizen of Illinois rather than Vermont. Second, the right to admission to practice law in state courts is not a privilege or immunity of \textit{United States} citizenship protected by Section 1 of the Fourteenth Amendment; states retain full authority to prescribe qualifications for admission to their own bars without federal constitutional interference.
  \item \textbf{Holding vs. Dicta:} The binding holding is Justice Miller's narrow federalism ruling that state bar admission does not depend on national citizenship and therefore falls outside the Fourteenth Amendment's Privileges or Immunities Clause under \textit{Slaughter-House}. Justice Miller's observation that admission to practice in \textit{federal} courts would relate to national citizenship (while noting that non-citizens had nonetheless been admitted) is dicta. Crucially, Justice Bradley's sweeping declarations about women's ``timidity and delicacy,'' coverture, and ``the law of the Creator'' appeared only in a three-Justice concurrence in the judgment, not in the opinion of the Court; although historically infamous, those gender-role pronouncements were never the Court's holding.
\end{itemize}

\section{Analysis of Opinions}

\subsection{Majority Opinion (Justice Miller)}

Writing for the same five-Justice majority that had decided the \textit{Slaughter-House Cases} the previous day, Justice Miller disposed of Bradwell's claims in a brief, strictly doctrinal opinion without relying on gender stereotypes. First, Miller rejected Bradwell's Article IV, \S{} 2 claim by pointing to her own affidavit: although born in Vermont, she had resided in Chicago for many years, and the first sentence of the Fourteenth Amendment makes U.S. citizens citizens of the state in which they reside. Because Article IV, \S{} 2 protects out-of-state citizens against discrimination by a host state, it did not shield an Illinois citizen from Illinois law.

Second, turning to the Fourteenth Amendment's Privileges or Immunities Clause, Miller agreed with Bradwell's counsel that certain privileges belong to U.S. citizens as such, but flatly denied that admission to a state bar is among them. Miller observed that bar admission had never depended on U.S. citizenship---noting that distinguished lawyers who were not citizens of the United States or any state had been admitted to both state and federal courts---and that insofar as bar admission related to citizenship at all, state-court admission related solely to state citizenship. Invoking \textit{Slaughter-House} as conclusive, Miller concluded that regulating state law licenses was a power never transferred to the federal government, leaving the Court with no authority to evaluate the reasonableness or propriety of Illinois's exclusion.

\subsection{Concurring Opinion (Justice Bradley)}

Justice Bradley, joined by Justices Swayne and Field, faced a doctrinal dilemma. The day before, all three had dissented in \textit{Slaughter-House}, arguing passionately that the Fourteenth Amendment's Privileges or Immunities Clause \textit{did} protect every citizen's fundamental right to pursue any lawful occupation free from arbitrary state monopolies. Because they rejected Miller's narrow reading of the Clause, they could not join his opinion and instead had to explain why male butchers in New Orleans possessed a constitutional right to pursue their trade while a qualified female lawyer in Chicago did not.

Bradley resolved this tension by asserting that the right to pursue any lawful employment had never historically been established as a fundamental privilege of women. Invoking civil law, English common law, and ``nature herself,'' Bradley declared that ``[m]an is, or should be, woman's protector and defender'' and that the ``natural and proper timidity and delicacy which belongs to the female sex evidently unfits it for many of the occupations of civil life.'' He cited the common-law doctrine of coverture---under which a married woman had no separate legal existence and could not make binding contracts without her husband's consent---as rendering married women incapable of discharging an attorney's fiduciary duties. Dismissing unmarried women as ``exceptions to the general rule,'' Bradley proclaimed that ``[t]he paramount destiny and mission of woman are to fulfill the noble and benign offices of wife and mother'' under ``the law of the Creator.'' Consequently, he concluded that excluding women from offices requiring ``decision and firmness'' fell squarely within the state's police power.

\subsection{Dissent (Chief Justice Chase)}

Chief Justice Salmon P. Chase, who had also dissented in \textit{Slaughter-House}, recorded a solitary dissent ``from the judgment of the Court and from all the opinions.'' Too ill to write an opinion, Chase left no written rationale. As noted by \href{https://en.wikipedia.org/wiki/Bradwell_v._Illinois}{historical accounts}, Chase was a lifelong abolitionist and supporter of women's suffrage; his lone dissent reflected a consistent view that the Fourteenth Amendment's broad guarantee of occupational liberty applied to women and men alike.

\section{Examples: Future Applications}

\begin{enumerate}
  \item \textbf{Same-Side Application (State Occupational Licensing under Privileges or Immunities):} A state enacts a licensing statute requiring hair braiders to complete 1,500 hours of cosmetology school unrelated to braiding, and an in-state resident challenges the law exclusively under the Fourteenth Amendment's Privileges or Immunities Clause. Under \textit{Bradwell} and \textit{Slaughter-House}, the Privileges or Immunities claim fails because the right to pursue a state-regulated occupation is not a privilege of national citizenship.
  \item \textbf{Same-Side Application (In-State Resident Invoking Article IV, \S{} 2):} A lifelong resident of Texas who was born in Vermont challenges a Texas medical-licensing restriction under Article IV, \S{} 2's Privileges and Immunities Clause. Under \textit{Bradwell}, Article IV, \S{} 2 is wholly inapplicable because the plaintiff is a citizen of the state whose law she challenges, regardless of her birthplace.
  \item \textbf{Opposite-Side Application (State Sex-Based Exclusion Challenged under Equal Protection):} A state-run nursing school or state bar excludes qualified applicants solely because of their sex, and an applicant sues under the Fourteenth Amendment's \textit{Equal Protection Clause}. Unlike Bradwell's Privileges or Immunities claim, modern Equal Protection doctrine applies intermediate scrutiny and invalidates sex classifications grounded in ``archaic and stereotypic notions'' about gender spheres, as the Supreme Court later held in \href{https://supreme.justia.com/cases/federal/us/458/718/}{\textit{Mississippi University for Women v. Hogan}}, 458 U.S. 718 (1982) (expressly citing \textit{Bradwell} as a historical example of invalid sex stereotyping).
  \item \textbf{Opposite-Side Application (Out-of-State Lawyer Excluded from State Bar):} A state bars all nonresidents from admission to its state bar, and a qualified attorney residing in a neighboring state challenges the residency requirement under Article IV, \S{} 2. Because the plaintiff is an out-of-state citizen---precisely the posture Bradwell lacked---the Supreme Court later held in \textit{Supreme Court of New Hampshire v. Piper}, 470 U.S. 274 (1985), that practicing law is a fundamental privilege protected against nonresident discrimination by Article IV, \S{} 2.
  \item \textbf{Fence-Sitter Application (Admission to Federal Courts vs. State Courts):} A federal district court adopts a local rule excluding a class of qualified U.S. citizens from its bar on a non-suspect basis. \textit{Bradwell} does not squarely resolve this case: Justice Miller's opinion explicitly distinguished state bars from federal bars, suggesting in dicta that admission to federal courts ``would relate to citizenship of the United States,'' yet also noted that non-citizens had been admitted to federal courts---leaving uncertain whether federal-bar admission is a protected incident of national citizenship.
\end{enumerate}

\section{Critique}

\textit{Bradwell v. Illinois} has been criticized from virtually every major school of constitutional thought, though for sharply different reasons.

From a \textbf{progressive and feminist perspective}, \textit{Bradwell} stands as an emblem of how nineteenth-century courts constitutionalized patriarchal ideology. Even though Justice Miller's majority opinion avoided overt sex stereotyping, its refusal to police state occupational exclusions left women without federal constitutional recourse against state-enforced economic subordination. Justice Bradley's concurrence openly elevated the Victorian ``separate spheres'' ideology and common-law coverture into ``the law of the Creator,'' ignoring the glaring contradiction posed by Myra Bradwell herself---a married woman who successfully edited and published the Midwest's premier legal newspaper and passed the Illinois bar examination. Moreover, because Bradwell's counsel did not press an Equal Protection argument, the Court never confronted whether sex-based exclusions violated the Equal Protection Clause---an omission that delayed meaningful constitutional scrutiny of sex discrimination for nearly a century until \textit{Reed v. Reed}, 404 U.S. 71 (1971), and \href{https://supreme.justia.com/cases/federal/us/458/718/}{\textit{Mississippi University for Women v. Hogan}}, 458 U.S. 718 (1982).

From an \textbf{originalist and textualist perspective}, \textit{Bradwell} exposes deep tensions within Reconstruction-era jurisprudence. On one hand, a broad scholarly consensus---represented by originalists such as Timothy Sandefur in the \href{https://www.law.nyu.edu/sites/default/files/ECM_PRO_065902.pdf}{\textit{NYU Journal of Law \& Liberty}} and Justice Clarence Thomas in \textit{McDonald v. City of Chicago}---argues that \textit{Slaughter-House} and \textit{Bradwell} mutilated the original public meaning of the Privileges or Immunities Clause, which the Thirty-Ninth Congress intended as the primary textual vehicle to protect fundamental civil and economic rights against state abridgment. On the other hand, originalists divide over whether restoring the Clause's original meaning would have changed \textit{Bradwell}'s outcome. Traditional originalists note that in 1868, ratifiers widely assumed common-law coverture and sex-based occupational exclusions remained permissible. By contrast, ``new originalists'' such as Steven Calabresi and Julia Rickert (discussed in \href{https://scholars.law.unlv.edu/cgi/viewcontent.cgi?article=1715&context=nlj}{UNLV Nevada Law Journal scholarship}) contend that the Fourteenth Amendment's original public meaning prohibited arbitrary caste legislation, and that once false nineteenth-century factual beliefs about women's intellectual capacities are discarded, original public meaning condemns \textit{Bradwell}'s result. Meanwhile, Bradley's concurrence is widely faulted by textualists for abandoning constitutional text in favor of unenumerated natural-law assumptions about gender roles---revealing the very subjectivity that plagued his economic-liberty dissent in \textit{Slaughter-House}.

\section{Key Quotations}

\begin{quote}
We agree with him that there are privileges and immunities belonging to citizens of the United States, in that relation and character, and that it is these and these alone which a state is forbidden to abridge. But the right to admission to practice in the courts of a state is not one of them. This right in no sense depends on citizenship of the United States.
\end{quote}
\noindent --- Justice Samuel F. Miller, writing for the Court, 83 U.S. at 139 (\href{https://supreme.justia.com/cases/federal/us/83/130/}{Justia}; \href{https://www.courtlistener.com/opinion/88662/bradwell-v-state/}{CourtListener}).

\begin{quote}
The opinion just delivered in the Slaughter-House Cases renders elaborate argument in the present case unnecessary, for, unless we are wholly and radically mistaken in the principles on which those cases are decided, the right to control and regulate the granting of license to practice law in the courts of a state is one of those powers which are not transferred for its protection to the federal government, and its exercise is in no manner governed or controlled by citizenship of the United States in the party seeking such license.
\end{quote}
\noindent --- Justice Samuel F. Miller, writing for the Court, 83 U.S. at 139 (\href{https://supreme.justia.com/cases/federal/us/83/130/}{Justia}).

\begin{quote}
Man is, or should be, woman's protector and defender. The natural and proper timidity and delicacy which belongs to the female sex evidently unfits it for many of the occupations of civil life. The Constitution of the family organization, which is founded in the divine ordinance as well as in the nature of things, indicates the domestic sphere as that which properly belongs to the domain and functions of womanhood.
\end{quote}
\noindent --- Justice Joseph P. Bradley, concurring in the judgment, 83 U.S. at 141 (\href{https://supreme.justia.com/cases/federal/us/83/130/}{Justia}).

\begin{quote}
The paramount destiny and mission of woman are to fulfill the noble and benign offices of wife and mother. This is the law of the Creator. And the rules of civil society must be adapted to the general constitution of things, and cannot be based upon exceptional cases.
\end{quote}
\noindent --- Justice Joseph P. Bradley, concurring in the judgment, 83 U.S. at 141--42 (\href{https://supreme.justia.com/cases/federal/us/83/130/}{Justia}).

\end{document}
